\scanpage{24}
donc le triangle \textcircled{a} ou \textcircled{a'} apparaît comme \uline{autotransposé}.

Appliquons le foncteur oubli, on trouve
\textcircled{c'}
\diagram{\begin{tikzcd}[row sep=large,column sep=large]
&L^\bullet\simeq\mathbb R\Gamma(W_0) \arrow[dl]\\
\mathbb R\operatorname{Hom}_{\Lambda}(K^{\bullet\mathrm{oub}},\Lambda)[-2n] \arrow[r] \arrow[d,equal] &K^{\bullet\mathrm{oub}}\simeq\mathbb R\Gamma(M) \arrow[u]\\
\mathbb R\Gamma_!(\mathring M)&
\end{tikzcd}}
\margin{\uline{dire que c'est encore un triangle autotransposé}}

Appliquons $\mathbb R\Gamma^\pi$, on trouve
\margin{[\uline{expliquer} comment ?]}
\textcircled{b'}
\margin{\deleted{\unsure{1}}}
\diagram{\begin{tikzcd}[row sep=large,column sep=large]
&L^\bullet\overset{\mathrm L}{\otimes}_{\Lambda}\mathbb R\Gamma(S^1)\simeq L^\bullet\times L^\bullet[-1]\simeq\unsure{I|l} \arrow[dl,"(\alpha',\beta')"']\\
\mathbb R\operatorname{Hom}_{\Lambda}(K^{\bullet\pi},\Lambda)[-2n-1] \arrow[r,very thick] \arrow[d,equal]&K^{\bullet\pi}\simeq\mathbb R\Gamma(P)\simeq\mathbb R\Gamma(W-W_0) \arrow[u,"(\alpha,\beta)",very thick]\\
\mathbb R\Gamma_!(\mathring P)\simeq\mathbb R\Gamma_!(W-W_0)&
\end{tikzcd}}
\margin{\fbox{\begin{minipage}{.33\linewidth}vérifier la commutation de $\mathbb R\Gamma^\pi$ à $\mathbb R\operatorname{Hom}_{\Lambda}(-,\Lambda)=(-)^\vee$\end{minipage}}}
\margin{(triangle autotransposé)}

\deleted{\unsure{Et les} \uline{triangles} \textcircled{d} et \textcircled{e} \uline{deviennent chacun}}

Comme nouvelle \deleted{\unsure{donnée}} \inserted{structure}, ce triangle introduit
\deleted{\diagram{\begin{tikzcd}
u_1:\ \mathbb R\operatorname{Hom}_{\Lambda}(K^\bullet,\Lambda)[-2n-1] \arrow[r] &\unsure{H|K}
\end{tikzcd}}}
\diagram{\begin{tikzcd}[column sep=large]
\boxed{u_1:\ \mathbb R\Gamma(W-W_0)^\vee[-2n-1]\mathbin{-}\mathbb R\Gamma(W-W_0)}
\end{tikzcd}}
et des accouplements
\diagram{\begin{tikzcd}
\boxed{\begin{gathered}
\lambda:\ \mathbb R\Gamma(W_0)\times\mathbb R\Gamma(W-W_0)\mathbin{-}\Lambda[-2n]\\
\mu:\ \mathbb R\Gamma(W_0)\times\mathbb R\Gamma(W-W_0)\mathbin{-}\Lambda[-(2n-1)]
\end{gathered}}
\end{tikzcd}}
\diagram{\begin{tikzcd}
\lambda:\ H^{\deleted{i}\inserted{2n-i}}(W_0)\times H^{\deleted{\unsure{2n-i}}\inserted{i}}(W-W_0) \arrow[r] &\Lambda\\
\mu:\ H^{\deleted{i}\inserted{2n-1-i}}(W_0)\times H^{\deleted{\unsure{2n-1-i}}\inserted{i}}(W-W_0) \arrow[r,dash] &\Lambda
\end{tikzcd}}
i.e.
\diagram{\begin{tikzcd}
\lambda':\ H^{\deleted{\unsure{2n-}}i}(W-W_0) \arrow[r]&H^{i-1}(W_0)&\text{je dis que c'est }\deleted{\unsure{le}}\ \beta\\
\mu':\ H^i(W-W_0) \arrow[r]&H^i(W_0)&\text{je dis que c'est }\alpha
\end{tikzcd}}
\mspage{23}
